\documentclass[11pt]{article}
\usepackage[utf8]{inputenc}
\usepackage{amsmath,amssymb}
\usepackage{hyperref}
\title{Scalable Sea Level Rise Projection for Large-Scale Settings}
\author{Anonymous}
\date{}
\begin{document}
\maketitle

\begin{abstract}
This paper studies Sea Level Rise Projection. Grounded in a real, citable literature \cite{ref1,ref2,ref3,ref4,ref5,ref6,ref7,ref8},
we propose a focused method and report \textbf{illustrative} experiments.
All references below are real and verifiable (OpenAlex/arXiv); the reported
numbers are simulated to illustrate intended behaviour.
\end{abstract}

\section{Introduction}
Sea Level Rise Projection is an active research area. This work targets a concrete gap identified from
the recent literature and contributes a method together with an evaluation protocol.

\section{Related Work (real literature)}
The following works, retrieved from public bibliographic APIs, frame our study:
\begin{itemize}
\item Neumann et al. (2015) — "Future Coastal Population Growth and Exposure to Sea-Level Rise and Coastal Flooding - A Global Assessment", PLoS ONE (cited 2840).
\item DeConto et al. (2016) — "Contribution of Antarctica to past and future sea-level rise", Nature (cited 2170).
\item Church et al. (2011) — "Sea-Level Rise from the Late 19th to the Early 21st Century", Surveys in Geophysics (cited 1622).
\item Church et al. (2006) — "A 20th century acceleration in global sea‐level rise", Geophysical Research Letters (cited 1619).
\item Knutson et al. (2019) — "Tropical Cyclones and Climate Change Assessment: Part II: Projected Response to Anthropogenic Warming", Bulletin of the American Meteorological Society (cited 1335).
\item Nerem et al. (2018) — "Climate-change–driven accelerated sea-level rise detected in the altimeter era", Proceedings of the National Academy of Sciences (cited 1136).
\end{itemize}

\section{Method}
We decompose the problem across shards with a communication-efficient aggregation rule that keeps overhead sublinear in scale. We instantiate this idea for Sea Level Rise Projection and analyse its cost/quality trade-off.

\section{Experiments (Illustrative)}
\begin{center}
\begin{tabular}{lcc}
System & Primary Metric & Efficiency \\ \hline
Strong baseline & 0.812 & 1.00$\times$ \\
Ablation & 0.828 & 0.92$\times$ \\
\textbf{Ours} & \textbf{0.851} & \textbf{0.71$\times$}
\end{tabular}
\end{center}
\emph{Metrics simulated to illustrate the intended effect; not measured.}

\section{Conclusion}
We connected a real literature to a concrete method for Sea Level Rise Projection. Future work will
validate the illustrative results with real experiments.

\begin{thebibliography}{9}
\bibitem{ref1} Barbara Neumann, Athanasios T. Vafeidis, Juliane Zimmermann et al.. Future Coastal Population Growth and Exposure to Sea-Level Rise and Coastal Flooding - A Global Assessment. \emph{PLoS ONE}, 2015. DOI:10.1371/journal.pone.0118571
\bibitem{ref2} Robert M. DeConto, David Pollard. Contribution of Antarctica to past and future sea-level rise. \emph{Nature}, 2016. DOI:10.1038/nature17145
\bibitem{ref3} John Church, Neil J. White. Sea-Level Rise from the Late 19th to the Early 21st Century. \emph{Surveys in Geophysics}, 2011. DOI:10.1007/s10712-011-9119-1
\bibitem{ref4} John Church, Neil J. White. A 20th century acceleration in global sea‐level rise. \emph{Geophysical Research Letters}, 2006. DOI:10.1029/2005gl024826
\bibitem{ref5} Thomas R. Knutson, Suzana J. Camargo, Johnny C. L. Chan et al.. Tropical Cyclones and Climate Change Assessment: Part II: Projected Response to Anthropogenic Warming. \emph{Bulletin of the American Meteorological Society}, 2019. DOI:10.1175/bams-d-18-0194.1
\bibitem{ref6} R. S. Nerem, B. D. Beckley, John Fasullo et al.. Climate-change–driven accelerated sea-level rise detected in the altimeter era. \emph{Proceedings of the National Academy of Sciences}, 2018. DOI:10.1073/pnas.1717312115
\bibitem{ref7} Mathieu Morlighem, Eric Rignot, Tobias Binder et al.. Deep glacial troughs and stabilizing ridges unveiled beneath the margins of the Antarctic ice sheet. \emph{Nature Geoscience}, 2019. DOI:10.1038/s41561-019-0510-8
\bibitem{ref8} David Pollard, Robert M. DeConto. Modelling West Antarctic ice sheet growth and collapse through the past five million years. \emph{Nature}, 2009. DOI:10.1038/nature07809
\end{thebibliography}
\end{document}
